\section*{6 A worst-case randomized reduction}
We now present a worst-case search to decision reduction for MFSP.
\setcounter{nttheorem}{40}
\begin{nttheorem}[Author Theorem 41]
There is a randomized algorithm solving Search-MFSP on inputs of length $N$ in time $O(2^{.67N})$ given access to an oracle that solves MFSP on instances of length $2N$.
\end{nttheorem}
\begin{proof}
We prove this theorem by giving an oracle algorithm solving DecompProblem and appealing to Theorem 20. We provide the pseudocode of our algorithm in Algorithm 3, which we recommend the reader look at before proceeding.
\begin{quote}\textbf{Algorithm 3. A randomized worst-case search to decision reduction for MFSP.}\par
\textbf{procedure WorstCaseOptimalSubcomputation($f$)}\par
Given the length-$N$ truth table of a function $f$ that takes $n$-inputs with $L(f)\geq 2$, this procedure returns an element $\{g,O,h\}$ of $\Opt(f)$. Set $s=\frac 23\cdot\frac{2^n}{\log n}$. Set $t=2^{2N/3}$.
\paragraph{Part 1: Try random formulas}
For $i=1,\ldots,t$ do:
\begin{enumerate}
\item Let $G_i$ be a uniformly random binary tree with $L(f)$-leaves. (Section 6.2 discusses how to sample $G_i$.)
\item Turn $G_i$ into a uniformly random formula $\phi_i$ by picking uniformly random gates from $\{\land,\lor\}$ and uniformly random input leaves from $\{0,1,x_1,\ldots,x_n,\neg x_1,\ldots,\neg x_n\}$.
\item If $\phi_i$ computes $f$ then: Write $\phi_i=\phi_{i,1}O\phi_{i,2}$. Let $g$ and $h$ be the function computed by $\phi_{i,1}$ and $\phi_{i,2}$ respectively. If $L(f)=L(g)+L(h)$ then return $\{g,O,h\}$. End if. End if.
\end{enumerate}
End for.
\paragraph{Part 2: Generate a small list of candidates for $g$}
Set $\operatorname{SmallFuncs}=\{g:g\text{ is a Boolean function with }n\text{-inputs and }L(g)\leq s\}$. Set $\Candidates=\{g\in\operatorname{SmallFuncs}:L(\Select[f,g])\leq L(f)+1\}$.
\paragraph{Part 3: Try to find a $g,h$ pair within Candidates}
For each pair of functions $(g,h)\in\Candidates$ and for each gate $O\in\{\land,\lor\}$ do: If $f=gOh$ and $L(f)=L(g)+L(h)$ then return $\{g,O,h\}$. End if. End for.
\paragraph{Part 4: Try to find a $g,h$ pair by looking at functions $h$ satisfying $f=gOh$}
Set $\operatorname{SmallCandidates}=\{g\in\Candidates:L(g)\leq L(f)-s\}$. For each function $g\in\operatorname{SmallCandidates}$ and for each $O\in\{\land,\lor\}$ do:
\begin{enumerate}
\item If $\forall x\in\{0,1\}^n\ \exists b\in\{0,1\}$ such that $g(x)Ob=f(x)$ then:
\begin{enumerate}
\item Let $h_{?,g}:\{0,1\}^n\to\{0,1,?\}$ be the unique partial function on $n$-inputs such that $\forall h$, $f=gOh\Longleftrightarrow h$ agrees with $h_{?,g}$.
\item For each total function $h$ that agrees with $h_{?,g}$ do: If $f=gOh$ and $L(f)=L(g)+L(h)$ then return $\{g,O,h\}$. End if. End for.
\end{enumerate}
End if.
\end{enumerate}
End for.\par
\textbf{end procedure}
\end{quote}
\subsection*{6.1 Correctness of Algorithm 3}
In this section, we prove that Algorithm 3 has the desired input/output behavior. In our analysis, we will use $s$ and $t$ as parameters which we will set to the optimal values (which are written in the pseudocode) in Section 6.2 where we do the running time analysis for Algorithm 3.
Fix some function $f$ on $n$-inputs with $L(f)\geq 2$. We analyze the algorithm in parts.
\paragraph{Part 1}
Since $\phi_i$ is chosen to have $L(f)$ leaves and the algorithm in part 1 checks if $\phi_i$ computes $f$ before returning any value, the following claim is clear.
\begin{ntclaim}[Author Claim 42]
Any output by Algorithm 3 returned in part 1 must be an element of $\Opt(f)$.
\end{ntclaim}
Moreover, we can lower bound the probability that Algorithm 3 returns a value in part 1 as follows. Recall that $\Canon{0}(f)$ is the set of optimal canonical formulas for $f$. We will show that part 1 succeeds if this set is large.
\begin{ntclaim}[Author Claim 43]
If $t\geq 5\cdot 2^{N(1+o(1))}/|\Canon{0}(f)|$, then part 1 of Algorithm 3 will return a value at least 99\% of time.
\end{ntclaim}
\begin{proof}
Since we are picking each $s$-leaf formula $\phi_i$ uniformly at random, the probability that any fixed formula computes $f$ is at least
\[\frac{|\Canon{0}(f)|}{\text{the total number of formulas with }L(f)\text{ leaves}}.\]
Combining Theorem 13 with Proposition 12 upper bounds the denominator by $2^{N(1+o(1))}$, so
\[\Pr[\phi_i\text{ computes }f]\geq\frac{|\Canon{0}(f)|}{2^{N(1+o(1))}}.\]
Since each of these $\phi_i$ are chosen independently, we have that
\begin{align*}\Pr[\exists i\in[t]\text{ such that }\phi_i\text{ computes }f]&\geq 1-\left(1-\frac{|\Canon{0}(f)|}{2^{N(1+o(1))}}\right)^t\\&\geq 1-e^{-t\cdot|\Canon{0}(f)|/2^{N(1+o(1))}}\\&\geq 1-e^{-5}\geq.99.\end{align*}
Hence, with probability at least 99\%, part 1 will find a $\phi_i$ computing $f$ at which point it will clearly return a value.
\end{proof}
\paragraph{Part 2}
In part 2, Algorithm 3 constructs the $\Candidates$ set. We prove two claims about this set. First, that it contains the functions we care about, and second that its size can be bounded using the size of the $\Canon{0}(f)$ set.
\begin{ntclaim}[Author Claim 44]
Suppose $g$ is in an optimal subcomputation for $f$. Then $L(g)\leq s\Longrightarrow g\in\Candidates$.
\end{ntclaim}
\begin{proof}
Since $L(g)\leq s$, we know that $g\in\operatorname{SmallFuncs}$. Next, since $g$ is in an optimal subcomputation for $f$, we have by Lemma 23 that $L(\Select[f,g])\leq L(f)+1$, so $g$ is an element of $\Candidates$.
\end{proof}
\begin{ntclaim}[Author Claim 45]
$|\Candidates|=O(|\Canon{0}(f)|\cdot N\log N)$.
\end{ntclaim}
\begin{proof}
By construction, we have that $\Candidates\subseteq\{g:L(\Select[f,g])=L(f)+1\}$. On the other hand, Lemma 24, we have that
\[|\{g:L(\Select[f,g])=L(f)+1\}|\leq O(|\Canon{0}(f)|\cdot N\log N).\]
\end{proof}
\paragraph{Part 3}
In Part 3, Algorithm 3 tries to find a $g,h$ pair by looking within the $\Candidates$ set. We show this works as long as there is a $\{g,O,h\}\in\Opt(f)$ where $L(g)$ and $L(h)$ are small.
First, we note that part 3 can only return correct answers.
\begin{ntclaim}[Author Claim 46]
Any output returned by Algorithm 3 in part 3 will be an element of $\Opt(f)$.
\end{ntclaim}
\begin{proof}
In order for a $\{g,O,h\}$ value to be returned in part 3 it must satisfy $f=gOh$ and $L(f)=L(g)+L(h)$. Thus, by Lemma 21, we know $\{g,O,h\}\in\Opt(f)$.
\end{proof}
Next, we give sufficient conditions on which part 3 will return an answer.
\begin{ntclaim}[Author Claim 47]
If there exists an element $\{g',O,h'\}$ of $\Opt(f)$ such that $\max\{L(g'),L(h')\}\leq s$, then Algorithm 3 will return a value in part 3 or before.
\end{ntclaim}
\begin{proof}
Suppose there exists an element $\{g',O,h'\}$ of $\Opt(f)$ such that $\max\{L(g'),L(h')\}\leq s$, and assume that this procedure has not returned a value before it reaches part 3. Then by Claim 44, we have that both $g'$ and $h'$ are in $\Candidates$. Moreover, since $\{g',O,h'\}\in\Opt(f)$, we know that $f=g'O'h'$ and $L(f)=L(g)+L(h)$ by Lemma 21. Hence it is clear there are value the for loop will return an output if it reaches $g=g'$, $h=h'$, and $O=O'$ (although it could return a value before that).
\end{proof}
\paragraph{Part 4}
In the final part of Algorithm 3, we look for matching $h$ functions for $g$ candidates with small complexity.
First, we note that any output returned by part 4 must be correct by essentially the same proof as Claim 46 in part 3.
\begin{ntclaim}[Author Claim 48]
Any output returned by Algorithm 3 in part 4 will be an element of $\Opt(f)$.
\end{ntclaim}
Next, we show sufficient conditions for part 4 returning an answer.
\begin{ntclaim}[Author Claim 49]
If $s\geq L(f)/2$ and there exists a $\{g',O,h'\}\in\Opt(f)$ such that $L(h')\geq s$, then Algorithm 3 will return a value in part 4 or earlier.
\end{ntclaim}
\begin{proof}
Using Lemma 21, we have $L(f)=L(g')+L(h')$. Thus, since $s\geq L(f)/2$ and $L(h')\geq s$, we have that
\[L(g')=L(f)-L(h')\leq L(f)-s\leq 2s-s=s.\]
Hence, by Claim 44, we have that $g'\in\Candidates$. Moreover, since $L(g')\leq L(f)-s$, it follows that $g'\in\operatorname{SmallCandidates}$. Thus, it is clear that part 4 will return a value if its for loop ever reaches $g=g'$, $O=O'$, and $h=h'$ (though it could return a value before that).
\end{proof}
\paragraph{In total}
Finally, we can prove the correctness of the input/output behavior of Algorithm 3.
\begin{ntclaim}[Author Claim 50]
If $s\geq L(f)/2$, then Algorithm 3 (run on input $f$) returns an element of $\Opt(f)$.
\end{ntclaim}
\begin{proof}
Put together, Claim 42, Claim 46, and Claim 48 ensures that any output returned by Algorithm 3 must be an element of $\Opt(f)$.
Hence, it suffices to show that Algorithm 3 will always output a value. We divide into two cases. Either there exists an element $\{g',O,h'\}\in\Opt(f)$ such that $\max\{L(g'),L(h')\}\leq s$ or not.
If there exists such an element, then Claim 47 ensures that Algorithm 3 will output a value.
Now suppose that for all $\{g',O,h'\}\in\Opt(f)$ we have $\max\{L(g'),L(h')\}\geq s$. Since $L(f)\geq 2$, we know $\Opt(f)$ is non-empty by Proposition 18. Hence we can fix some $\{g',O,h'\}\in\Opt(f)$ satisfying $\max\{L(g'),L(h')\}\geq s$. Without loss of generality we can assume that $L(g')\leq L(h')$. Thus, we have that $L(h')\geq s$, and by hypothesis $s\geq L(f)/2$, so Claim 49 ensures Algorithm 3 outputs a value.
\end{proof}
Thus Algorithm 3 is correct as long as $s\geq L(f)/2$. Indeed, in the next section we will set $s$ so that $s\geq\max\{L(f)/2:f\text{ takes }n\text{-inputs}\}$.
